\documentclass[11pt,a4paper,oneside,openany]{book}
\usepackage{biblioflow-ista}

\renewcommand{\titolodocumento}{Report Finale del Progetto}
\renewcommand{\sottotitolodocumento}{Report Finale del Progetto}

\begin{document}
\frontmatter
\copertina
\tableofcontents

\mainmatter
\chapter{Introduzione}

In questo documento viene illustrato il percorso evolutivo di BiblioFlow,
descrivendo le attività di manutenzione ed estensione svolte a partire
dall'applicazione esistente. Il documento sintetizza le fasi di analisi
preventiva, l'implementazione delle richieste di modifica (Change Requests)
e la successiva validazione tramite test, evidenziando come la qualità del
software sia stata preservata e migliorata.
\begin{itemize}
  \item \textbf{Affidabilità della Prenotazione (CR01):} l'identità di chi
        prenota deriva sempre dalla sessione e la creazione è atomica: un
        vincolo del database rende impossibile assegnare lo stesso posto a
        due studenti nella stessa fascia, anche con richieste simultanee.
  \item \textbf{Posti che non restano vuoti (CR02):} chi trova il posto
        occupato può mettersi in lista d'attesa; quando il posto si libera per
        cancellazione o mancata presenza, il primo in coda lo ottiene
        automaticamente e ne viene avvisato.
\end{itemize}

\chapter{Sviluppo del Progetto}

BiblioFlow si configura come una piattaforma web per la gestione dei posti
studio e dei servizi di una biblioteca universitaria. Il sistema unisce la
prenotazione guidata dei posti, il check-in con QR code e le automazioni che
liberano i posti non occupati, con un'attenzione particolare agli studenti
pendolari e all'accessibilità.

\section{Componenti Core (Interfaccia, API e Automazioni)}

La logica applicativa risiede in un'applicazione Next.js organizzata in
package distinti: pagine e route handler in \file{src/app}, servizi condivisi
in \file{src/lib}, componenti dell'interfaccia in \file{src/components} e
modello dei dati in \file{prisma}. Le Change Requests hanno introdotto un
servizio di dominio dedicato alle prenotazioni e alla lista d'attesa, al
quale le route si limitano a delegare: le nuove regole sono state aggiunte
come moduli dedicati, minimizzando l'impatto sul codice esistente.

Le automazioni periodiche sono separate dall'interfaccia e richiamate da
un'attività pianificata tramite un endpoint protetto. Questa separazione ha
facilitato l'implementazione della CR02, permettendo di collegare la
promozione dalla coda ai tre punti in cui un posto si libera senza
riprogettarli.

\section{Change Requests}

Questa sezione riassume le Change Requests (CR) utili all'evoluzione di
BiblioFlow. Le modifiche hanno l'obiettivo di rendere affidabile la
prenotazione dei posti e di non lasciare vuoti i posti che si liberano.

{\small
\begin{xltabular}{\linewidth}{L{1cm}YL{2.6cm}L{1.8cm}}
  \toprule
  \intestazione \tb{ID} & \tb{Titolo} & \tb{Categoria} & \tb{Esito} \\
  \midrule
  CR01 & Prenotazione sicura e atomica dei posti studio & Corrective,
  Preventive & Rilasciata \\
  CR02 & Lista d'attesa con riassegnazione automatica dei posti &
  Perfective & Rilasciata \\
  \bottomrule
\end{xltabular}}

\paragraph{CR01 -- Prenotazione sicura e atomica dei posti studio}
\textbf{Obiettivo:} ricavare identità e titolarità sempre dalla sessione
autenticata, centralizzare le regole sugli orari e rendere atomica la
creazione delle prenotazioni, così che due prenotazioni attive non possano
mai sovrapporsi sullo stesso posto.

\paragraph{CR02 -- Lista d'attesa con riassegnazione automatica dei posti}
\textbf{Obiettivo:} permettere allo studente di entrare in coda per un posto
occupato e promuovere automaticamente il primo in coda quando il posto si
libera, con notifica, scadenza della promozione non confermata e
tracciabilità di ogni passaggio.

\section{Strategia di Testing}

La verifica del sistema ha seguito un approccio incrementale definito nel
Master Test Plan, coprendo sia le nuove funzionalità sia la non-regressione
dell'esistente.

\subsection{Gestione degli Incidenti (Pre-Modification)}

Durante la fase preliminare di testing sono emersi problemi nella suite
stessa, non nel sistema: la versione di partenza non conteneva alcuna
infrastruttura di test, una suite veniva saltata in silenzio a causa del
caricamento di un pacchetto e un caso di caratterizzazione fotografava un
difetto invece di un requisito. L'infrastruttura è stata aggiunta senza
toccare i flussi verificati e i test sono stati corretti, documentando
ciascun intervento nel Test Incident Report.

\subsection{Testing delle Nuove Funzionalità (Post-Modification)}

\textbf{Prenotazione (CR01):} sono stati eseguiti test di unità sul servizio
di dominio e sulle funzioni di autorizzazione, e test di integrazione su
PostgreSQL reale per verificare che, con richieste simultanee sullo stesso
posto, una sola prenotazione venga creata e nessun dato parziale resti nel
database.

\textbf{Lista d'attesa (CR02):} sono stati verificati su database reale
l'ordine della coda, la promozione dopo cancellazione, annullamento dello
staff e no-show, la scadenza delle promozioni e l'assenza di duplicati con
esecuzioni concorrenti. È stata introdotta una suite End-to-End con
Playwright che percorre nel browser prenotazione, conflitto, ingresso in
coda, promozione e uscita dalla coda.

\textbf{Qualità della suite:} l'estensione della copertura alle route,
prima escluse dalla misura, ha portato il codice server al 95,2\% delle
istruzioni e all'89,4\% dei rami, con 1013 test superati su 1013. I nuovi test
hanno fatto emergere quindici difetti, corretti nella manutenzione
successiva alle Change Requests.

\subsection{Esito della Regressione}

La riesecuzione della suite storica alla chiusura delle Change Requests ha
confermato che le modifiche non hanno introdotto difetti: i 20 casi
classificati come invariati hanno conservato l'esito della versione di
partenza. Sulla versione finale 12 casi restano invariati e superati; gli
altri sono stati modificati o eliminati da interventi correttivi documentati,
ciascuno verificato da un nuovo test. Tutti i sei criteri di accettazione
risultano soddisfatti.

\end{document}
